\documentclass{article}
\usepackage{../math_template}

\author{Jonathan Kasongo}
\title{Moldova TST 2023 --- P4}

\begin{document}
\maketitle

\begin{problem}{Moldova TST 2023 --- P4}
Polynomials $\left(P_n(X)\right)_{n \in \mathbb{N}}$ are defined as:
$$
\begin{aligned}
& P_0(X)=0, \quad P_1(X)=X+2 \\
& P_n(X)=P_{n-1}(X)+3 P_{n-1}(X) \cdot P_{n-2}(X)+P_{n-2}(X), \quad
\forall n \geq 2 .
\end{aligned}
$$

Show that if $k$ divides $m$ then $P_k(X)$ divides $P_m(X)$.
\end{problem}

\begin{solution}{Solution}
One can apply Simon's favourite factoring trick to see that
$
3P_n + 1 = (3P_{n-1} + 1)(3P_{n-2} + 1)
$. If we let $Q_n = \ln (3P_n + 1)$ then we have
$Q_n = Q_{n-1} + Q_{n-2}$. Solving the recurrence relation with any
standard method yields
$3P_n + 1 = (3x+7)^{f_n}$ where $f_n$ is the $n$-th fibonacci number.
Thus

$$
P_n (x) = \frac13 [(3x+7)^{f_n} - 1], \quad n \geq 2
$$

and so the roots of $P_n$ are of the form $\frac13 (\omega_k - 7)$, where
$\omega_k$ are the $f_n$-th roots of unity. If $a \mid b$ then
the $a$-th roots of unity are a proper subset of the $b$-th roots of unity,
so we will be done if we show that $a \mid b \implies f_a \mid f_b$.
\\

\textit{Claim 1: For non-negative integers $k$ and $m$,
$f_{k+m} \equiv f_m f_{k+1} \ \text{(mod } f_k)$.} \vspace{0.2cm}

\textit{Proof:} We will induct on $m$. Throughout this proof any
equivalences are assumed to be modulo $f_k$.
If $m=1$ then $f_{k+1} \equiv f_1 f_{k+1}$ and if $m=2$ then $f_{k+2} \equiv
f_{k+1} + f_k \equiv f_2 f_{k+1}$. This completes the base cases.
\vspace{0.2cm}

Now suppose the result holds for $m=a$ and $m=a+1$. Consider
$f_{k+a+2} \equiv f_{k+a+1} + f_{k+a} \equiv f_{k+1}(f_{a+1} + f_a)
\equiv f_{k+1}f_{a+2}$. This finishes the proof by induction. $\square$
\\

As a corollary of claim 1, we find that $f_{a + \lambda b}
\equiv f_{b} f_{a+1}^{\lambda} \ \text{(mod } f_a)$ for $a,b,\lambda, \in
\mathbb{Z}_{\geq0}$. Now letting $a=b$ shows the desired result
$x \mid y \implies f_x\mid f_y$. $\blacksquare$
\end{solution}
\end{document}
